\SourceLeaf{426}%
die Fläche der Werth von $V$ sein Zeichen ändern, so dass, so lange die Stetigkeit nicht unterbrochen wird, auf der einen Seite die positiven, auf der andern die negativen [Werthe] stattfinden.

Die Richtung der Normale auf die Fläche nach derjenigen Seite zu, die wir als obere betrachten, werde auf der Hülfskugel durch den Punkt $L$ bezeichnet und es sei
\[
\cos L(1)=X,\qquad\cos L(2)=Y,\qquad\cos L(3)=Z.
\]
Es sei ferner $ds$ eine unendlich kleine Linie auf der Fläche, und indem deren Richtung auf der Hülfskugel durch den Punkt $\lambda$ bezeichnet wird, sei
\[
\cos\lambda(1)=\xi,\qquad\cos\lambda(2)=\eta,\qquad\cos\lambda(3)=\zeta.
\]
Wir haben dann
\[
dx=\xi ds,\qquad dy=\eta ds,\qquad dz=\zeta ds,
\]
also
\[
P\xi+Q\eta+R\zeta=0
\]
und, da $\lambda L=90^{0}$ sein muss, auch
\[
X\xi+Y\eta+Z\zeta=0.
\]
Da $P,Q,R,X,Y,Z$ bloss von der Stelle der Fläche abhängen, von wo man das Element ausgehen lässt, und diese Gleichungen für jede Richtung des Elements in der Fläche richtig sind, so sieht man leicht, dass $P,Q,R$ den $X,Y,Z$ proportional sein müssen, also
\[
P=X\mu,\qquad Q=Y\mu,\qquad R=Z\mu,
\]
also, da
\[
XX+YY+ZZ=1
\]
wird:
\[
\mu=PX+QY+RZ
\]
und
\[
\mu\mu=PP+QQ+RR
\]
oder
\[
\mu=\pm\BareRoot{(PP+QQ+RR)}.
\]

Entfernt man sich von der Fläche in der Richtung nach der Normale um
